\documentclass[book.tex]{subfiles}
\begin{document}

\chapter{Optimization: Golden Search}
\label{chap:golden_search}
\noindent\rule{11cm}{0.4pt} \\
\textbf{Prerequisites:} \autoref{chap:optimization},  \autoref{chap:bracketing}, \autoref{chap:newton_raphson} \\
\textbf{Difficulty Level:} * \\
\noindent\rule{11cm}{0.4pt}
\vspace{5mm}

In this chapter we will learn about basic optimization and the difference between one-dimensional and multidimensional optimization. We will study the difference between global and local optima and then define the golden ratio and how to use it to make one dimensional optimization efficient. We will then explain how to locate the optimum of a single variable function with golden section search.

\section{Optimization}
Optimization is the process of creating something that is as effective as possible. Optimization deals with finding the maxima and minima of a function that depends on one or more variables. The goal is to determine the values of the variables that yield maxima or minima for the function which can then be substituted back into the function to compute its optimal values.
\begin{marginfigure}[-6\baselineskip]
	\centering
	\includegraphics[width = 2.4in]{figures/part1b/golden_section/optimization.png}
	\centering
	\caption{Difference between roots and optima}
	\label{fig:1}
\end{marginfigure}

Optimization is similar in spirit to the root-location methods we covered in previous chapters, both involve guessing and searching for a point on a function. The fundamental difference between the two types of problems is illustrated in Figure \ref{fig:1}. Root location involves searching for the location where the function equals zero. In contrast, optimization involves searching for the function's extreme points.

As can be seen in Figure \ref{fig:1}, the optima are the points where the curve is flat. In mathematical terms, this corresponds to the $x$ value where the derivative $f' (x)$ is equal to zero. Additionally, the second derivative, $f''(x)$, indicates whether the optimum is a minimum or a maximum: if $f(x) < 0$, the point is a maximum; if $f(x) > 0$, the point is a minimum. Differentiating the function and locating the root of the new function $f'(x) = 0$ is a possible strategy for finding the optima. A variety of direct numerical optimization methods are available for both one-dimensional and multidimensional problems.

\begin{marginfigure}
	\centering
	\includegraphics[width = 2.4in]{figures/part1b/golden_section/onedimensional_1.png}
	\centering
	\caption{(a) One-dimensional Optimization (b) Two-dimensional Optimization }
	\label{fig:2}
\end{marginfigure}

One-dimensional problems involve functions that depend on a single dependent variable. As in Figure \ref{fig:2}a, the search then consists of climbing or descending one-dimensional peaks and valleys. Multidimensional problems involve functions that depend on two or more dependent variables. Two-dimensional optimization can again be visualized as searching out peaks and valleys (Figure \ref{fig:2}b). The process of finding a maximum versus finding a minimum is essentially identical because the same value $x^*$ both minimizes $f(x)$ and maximizes $-f(x)$. This equivalence is illustrated graphically for a one-dimensional function in Figure \ref{fig:2}a.

\section{One-dimensional Optimization}

Consider the function depicted in Figure \ref{fig:3}. Recall that root location was complicated by the fact that several roots can occur for a single function. Similarly, both local and global optima can occur in optimization. As can be seen from the Figure \ref{fig:3}, the \textit{global optimum} represents the very best solution. A \textit{local optimum}, though not the very best, is better than its immediate neighbors. Cases that include local optima are called \textit{multimodal}. For simple functions, we will almost always be interested in finding the global optimum, but we must be concerned about mistaking a local result for the global optimum. Just as in root location, optimization in one dimension can be divided into bracketing and open methods.

\begin{marginfigure}
	\centering
	\includegraphics[width = 2.4in]{figures/part1b/golden_section/onedimensional_2.png}
	\centering
	\caption{Function asymptotically approaching zero at plus and minus $\infty$	}
	\label{fig:3}
\end{marginfigure}

\section{Golden-Section Search}

Consider a continuous function $f(x)$ over the interval $[x_l,x_u]$ that contains a single minimum as shown in Figure \ref{fig:4}a , and hence is called \textit{unimodal}. Instead of using a single intermediate value to detect a sign change, we would need two intermediate function values to detect whether a minimum occurred. Consider two intermediate points ($x_1$ and $x_2$) within the interval at a distance "d" from $x_l$ and $x_u$ respectively. The function is evaluated at these two points. Two results can occur:
\begin{itemize}
	\item If, as in Figure \ref{fig:4}a, $f(x_1)< f(x_2)$, then $f(x_1)$ is the minimum, and the domain of $x$ to the left of $x_2$, from $x_l$ to $x_2$, can be eliminated because it does not contain the minimum. For this case, $x_2$ becomes the new $x_l$ for the next round.
	\item If $f(x_2)< f(x_1)$, then $f(x_2)$ is the minimum and the domain of $x$ to the right of $x_1$, from $x_1$ to $x_u$ would be eliminated. For this case, $x_1$ becomes the new $x_u$ for the next round. 
\end{itemize}
 \begin{marginfigure}[+8cm]
	\centering
	\includegraphics[width = 2.4in]{figures/part1b/golden_section/goldensection_2.png}
	\centering
	\caption{(a) Initial step (b) Second step of the search algorithm}
	\label{fig:4}
\end{marginfigure}
As the iterations are repeated, the interval containing the extremum is reduced rapidly. The key to making this approach efficient is the wise choice of the intermediate points. As in bisection, the goal is to minimize function evaluations by replacing old values with new values.\\
%Imagine that the original range in Figure \ref{fig:5} is of unit length. Let's choose intermediate points in such a way that $x_1 - x_l = x_u - x_2 = 1/3$. Also let $f(x_1) < f(x_2)$, such that the extremum will lie in the interval $[x_l,x_2]$. For the next iteration since $x_1$ is already in the new uncertainty interval and as the value of the function at $x_1$ i.e. $f(x_1)$ is already evaluated, we can coincide $x_1$ with the new 
%$x_{2new}$. Thus, for this interval we end up evaluating the function only at $x_{1new}$.

 Now, let's define $\phi$ as indicated in Figure \ref{fig:5} where $x_1 - x_l = x_u - x_2 = \phi(x_u - x_l)$. The intermediate points are chosen such that the reduction in the range is symmetric. Now consider the function with intermediate points such that $f(x_2) < f(x_1)$ and hence the extremum will lie in the interval $[x_l,x_1]$. Coinciding $x_2$ with the new $x_{2new}$ for which the function is already evaluated, we reduce the number of function evaluations to just one and at $x_{1new}$. Assuming the range to be of unit length in Figure \ref{fig:5}, it can be seen
   

\begin{align}
(x_1 - x_l) &= \phi\\
(x_u - x_2) &= \phi
\end{align}
Hence $(x_2 - x_l) = 1 - \phi$. For the second iteration, the bracket $[x_l,x_1]$ is further divided in the ratio of $\phi$ to obtain intermediate points as $x_{1new}$ and $x_2$ where
\begin{align}
(x_2 - x_l) = \phi
\end{align}
Therefore 
\begin{align}
\phi(x_1 - x_l) &= (x_2 - x_l)\\
\phi(x_1 - x_l) &= 1 - \phi\\
\phi(\phi) &= 1 - \phi\\
(\phi)^2 + \phi - 1 &= 0\\
\phi &= \dfrac{-1+\sqrt5}{2}
\end{align}
 \begin{marginfigure}
	\centering
	\includegraphics[width = 2.4in]{figures/part1b/golden_section/goldensection_3.png}
	\centering
	\caption{Algorithm of Golden Section search }
	\label{fig:5}
\end{marginfigure}
The other root is neglected as we need $\phi$ to be greater than $1/2$ as the reduction in range needs to be symmetric. Therefore, \textbf{$\phi = 0.61803$}. It can be shown that 
\begin{align}
\dfrac{\phi}{1 - \phi} = \dfrac{1 - \phi}{\phi}
\end{align}
Dividing a range in the ratio of $\phi$ to $1 - \phi$ has the effect that the ratio of the shorter segment to the longer equals to the ratio of the longer to the sum of the two. This rule is called \textbf{golden section}. The uncertainty range is reduced by the ratio $\phi = 0.61803$ at every stage. Hence, N steps of reduction using the golden section method reduces the range by the factor $(\phi)^N \approx (0.61803)^N$.

 For the golden-section search, the two intermediate points are chosen according to the golden ratio:
 \begin{align}
  x_1 = x_l + \phi(x_u - x_l)\\
 x_2 = x_u - \phi(x_u - x_l)\\
 \end{align}
 The function is evaluated at these two interior points. For the next iteration we need only determine the $x_{1new}$. This is done based on the new values of $x_l$ and $x_u$. A similar approach would be used for the alternate case where the $x_{2new}$ would be computed using $x_u$.
 
An upper bound for golden-section search can be derived as follows: Once an iteration is complete, the optimum will either fall in one of two intervals. If the optimum function value is at $x_2$, it will be in the lower interval ($x_l, x_2, x_1$). If the optimum function value is at $x_1$, it will be in the upper interval ($x_2, x_1, x_u$). Because the interior points are symmetrical, either case can be used to define the error.

Considering the upper interval ($x_2, x_1, x_u$), if the true value were at the far left, the maximum distance from the estimate can be shown to be $0.2361(x_u - x_l)$. Similarly, if the true value were at the far right, the maximum distance from the estimate would be $0.382 (x_u - x_l)$. Therefore, this case would represent the maximum error. This result can then be normalized to the optimal value for that iteration $x_{opt}$ to yield

\begin{align}
\varepsilon_a = \phi\arrowvert\dfrac{x_u - x_l}{x_{opt}}\arrowvert \times 100\%
\end{align}
This estimate provides a basis for terminating the iterations.

\begin{example}
Let's consider an unimodal function and implement the golden-section search method.

\begin{lstlisting}[language=python]
def f(x: $\mathbb{R}$) -> $\mathbb{R}$:
  return exp(x)-2*x                                        
xl = -10                                                   
xu = 10 
eps = 1e-6                                                 
phi = (-1+sqrt(5))/2
x1 = xl + phi*(xu-xl)                                      
x2 = xu - phi*(xu-xl)                                      
fx1 = f(x1)                                                
fx2 = f(x2)
while True:
  if fx1 < fx2:
    xl = x2                                                
    x2 = x1                                                
    fx2 = fx1                                              
    x1 = xl + phi*(xu-xl)                                  
    fx1 = f(x1)                                            
  else:
    xu = x1
    x1 = x2 
    fx1 = fx2
    x2 = xu - phi*(xu-xl)                                  
    fx2 = f(x2)                                            
  if abs(xu-xl) < eps:                                     
    break 
print('Extremum bracket')                                  
print([xl, xu])                                            
\end{lstlisting}
\end{example}

\section{Exercises}
\begin{enumerate}
    \item TODO
\end{enumerate}
\end{document}